\begin{frame}{왜 출력층 오차가 ``$a_2 - y$·$\sigma'$''로 깔끔한가}
  \small
  \begin{columns}[T]
    \begin{column}{0.49\textwidth}
      \kb{MSE} $L = \tfrac{1}{2}(a_2 - y)^2$ 일 때:
      \[
      \delta_2 = (a_2 - y)\,\sigma'(z_2)
      = (a_2 - y)\,a_2(1 - a_2)
      \]
      시그모이드 미분항이 \textbf{그대로 곱}된다.
    \end{column}
    \begin{column}{0.51\textwidth}
      \kb{교차 엔트로피} $L = -[y\ln a_2 + (1-y)\ln(1-a_2)]$ 일 때:
      \[
      \frac{\partial L}{\partial a_2} = -\frac{y}{a_2} + \frac{1-y}{1-a_2}
      \]
      $\sigma'(z_2) = a_2(1-a_2)$와 곱하면 \textbf{정확히 약분}.
    \end{column}
  \end{columns}
  \vskip0.3em
  \begin{block}{``교차 엔트로피 + 시그모이드 출력''이 표준인 계산적 이유}
    $\delta_2 = a_2 - y$ — 포화 항(최대 0.25)이 붙지 않아,
    예측이 0이나 1에 포화돼도 그래디언트가 상대적으로 크게 유지.
  \end{block}
  {\footnotesize 은닉층에서는 약분이 안 일어나 $\delta_1$에는 여전히
  $\sigma'(z_1)$가 곱된다 — 9.2절 소실 문제의 시작점.}
\end{frame}
